%============================================================
%  ch04_minsubspaces.tex
%  Chapter 4: Minimal Subspaces in Tensor Representations
%
%  Tensor Formats in Banach Spaces
%  Antonio Falcó, CEU Universities
%
%  Based on: Falcó & Hackbusch, Found. Comput. Math. 12 (2012)
%  Estimated length: ~40 pp.
%============================================================

\chapter{Minimal Subspaces in Tensor Representations}
\label{chap:minsubspaces}

In Chapter~\ref{chap:algtensor} we introduced the Tucker format and
defined the minimal subspace $\Umin{\alpha}{\bv}$ informally, taking
its existence and properties for granted. This chapter provides the
complete theory: existence and uniqueness (Section~\ref{sec:min:existence}),
the lattice structure and inclusion under refinement of partitions
(Section~\ref{sec:min:lattice}), admissible ranks and the disjoint
decomposition of $\VD$ (Section~\ref{sec:min:admissible}), and the
existence of best approximations in Banach tensor spaces
(Section~\ref{sec:min:bestapprox}). The results in this chapter are
taken from~\cite{FH2012}, expanded with complete proofs and additional
examples.

%------------------------------------------------------------
\section{Existence and uniqueness of minimal subspaces}
\label{sec:min:existence}
%------------------------------------------------------------

\subsection*{The case $d = 2$: Schmidt decomposition}

The case of two factors $\VD = V_1 \atens V_2$ is classical and serves
as the inductive base. The key observation is a lattice property of
tensor subspaces.

\begin{lemma}[Intersection lemma]
\label{lem:intersection_d2}
Let $X_i, Y_i$ be subspaces of $V_i$ for $i = 1, 2$. Then
\begin{equation}
\label{eq:intersection_d2}
    (X_1 \atens X_2) \cap (Y_1 \atens Y_2)
    = (X_1 \cap Y_1) \atens (X_2 \cap Y_2).
\end{equation}
\end{lemma}

\begin{proof}
The inclusion $(X_1 \cap Y_1) \atens (X_2 \cap Y_2) \subseteq
(X_1 \atens X_2) \cap (Y_1 \atens Y_2)$ is immediate. For the
reverse inclusion, let $\bv \in (X_1 \atens X_2) \cap (Y_1 \atens Y_2)$
with representations
\begin{equation*}
    \bv = \sum_{\nu=1}^{n_x} x^{(1)}_\nu \otimes x^{(2)}_\nu
       = \sum_{\nu=1}^{n_y} y^{(1)}_\nu \otimes y^{(2)}_\nu,
    \quad
    x^{(i)}_\nu \in X_i,\; y^{(i)}_\nu \in Y_i.
\end{equation*}
By adding terms with zero coefficient, assume $n_x = n_y = n$ and
that $\{x^{(1)}_\nu\}$ and $\{y^{(1)}_\nu\}$ are each linearly
independent (passing to a minimal representation). The two
representations yield the same matrix $M_1(\bv)$ under any identification
with a bilinear form, so $\spann\{x^{(1)}_\nu\} = \spann\{y^{(1)}_\nu\}
\subset X_1 \cap Y_1$. A symmetric argument for the second factor gives
$\spann\{x^{(2)}_\nu\} \subset X_2 \cap Y_2$, so
$\bv \in (X_1 \cap Y_1) \atens (X_2 \cap Y_2)$.
\end{proof}

\begin{proposition}[Minimal subspace: case $d=2$]
\label{prop:min_d2}
Let $\bv \in V_1 \atens V_2$. There exist unique subspaces
$U_i \subset V_i$ of finite dimension such that:
\begin{enumerate}
\item[\textup{(i)}] $\bv \in U_1 \atens U_2$.
\item[\textup{(ii)}] If $\bv \in X_1 \atens X_2$ for any subspaces $X_i$,
    then $U_i \subset X_i$.
\end{enumerate}
Moreover, $\dim U_1 = \dim U_2$.
\end{proposition}

\begin{proof}
\textbf{Existence.} Fix any representation $\bv = \sum_\nu v^{(1)}_\nu
\otimes v^{(2)}_\nu$ and set $U_i = \spann\{v^{(i)}_\nu\}$. By passing
to a representation with linearly independent $v^{(1)}_\nu$, one can also
take $v^{(2)}_\nu$ linearly independent (otherwise row reduce).
Then $\dim U_1 = \dim U_2$ (both equal the rank of the bilinear form).

\textbf{Minimality.} If $\bv \in X_1 \atens X_2$, then by the
intersection lemma, $\bv \in (U_1 \cap X_1) \atens (U_2 \cap X_2)$,
which forces $U_i \subset X_i$ by minimality of $U_i$.

\textbf{Uniqueness.} If $U_i$ and $U_i'$ both satisfy (i)--(ii),
minimality of $U_i$ applied to $U_i'$ gives $U_i \subset U_i'$,
and vice versa.
\end{proof}

\subsection*{The general case: existence and uniqueness}

\begin{theorem}[Existence and uniqueness of minimal subspaces]
\label{thm:min_existence}
Let $D = \{1, \ldots, d\}$, $\VD = {}_a\bigotimes_{\alpha \in D}
V_\alpha$, and $\bv \in \VD$. For each $\alpha \in D$, under the
identification $\VD \cong V_\alpha \atens \Vbigalpha^{[\alpha]}$,
there exists a unique subspace $\Umin{\alpha}{\bv} \subset V_\alpha$
of finite dimension $r_\alpha < \infty$ such that:
\begin{enumerate}
\item[\textup{(a)}] $\bv \in \Umin{\alpha}{\bv} \atens
    \Vbigalpha^{[\alpha]}$.
\item[\textup{(b)}] If $\bv \in U_\alpha \atens \Vbigalpha^{[\alpha]}$
    for any subspace $U_\alpha \subset V_\alpha$, then
    $\Umin{\alpha}{\bv} \subset U_\alpha$.
\end{enumerate}
Furthermore, there exists a unique subspace
$\Umin{D\setminus\{\alpha\}}{\bv} \subset \Vbigalpha^{[\alpha]}$
of dimension $r_\alpha$ such that
$\bv \in \Umin{\alpha}{\bv} \atens \Umin{D\setminus\{\alpha\}}{\bv}$.
\end{theorem}

\begin{proof}
Apply Proposition~\ref{prop:min_d2} to the two-factor product
$V_\alpha \atens \Vbigalpha^{[\alpha]}$. The subspace $\Umin{\alpha}{\bv}$
is the minimal subspace in the $V_\alpha$-factor. Uniqueness follows
from the same argument as in $d=2$.
\end{proof}

\begin{definition}
\label{def:minimal_subspace}
The subspace $\Umin{\alpha}{\bv}$ of Theorem~\ref{thm:min_existence}
is called the \emph{minimal subspace} of $\bv$ in direction $\alpha$.
The integer $r_\alpha(\bv) \coloneqq \dim \Umin{\alpha}{\bv}$ is the
\emph{$\alpha$-rank} of $\bv$.
\end{definition}

\subsection*{Characterisation via contractions}

The following characterisation, which is central to both the proof of
Chapter~\ref{chap:tucker_manifold} and the tree-based theory of
Chapter~\ref{chap:algTBT}, expresses $\Umin{\alpha}{\bv}$ as the
range of a family of contractions.

\begin{theorem}[Contraction characterisation]
\label{thm:min_contraction}
Let $\bv \in \VD$ and $\alpha \in D$. Then
\begin{equation}
\label{eq:min_contraction}
    \Umin{\alpha}{\bv}
    = \spann\bigl\{
        (\Id{V_\alpha} \otimes \varphi^{[\alpha]})(\bv) :
        \varphi^{[\alpha]} \in (\Vbigalpha^{[\alpha]})^*
      \bigr\},
\end{equation}
where $(\Id{V_\alpha} \otimes \varphi^{[\alpha]})(\bv) \in V_\alpha$
is the contraction of $\bv$ with the functional $\varphi^{[\alpha]}$
in all directions except $\alpha$.
\end{theorem}

\begin{proof}
Write $\bv = \sum_\nu u_\nu \otimes \bU_\nu$ with $u_\nu \in V_\alpha$
and $\bU_\nu \in \Vbigalpha^{[\alpha]}$. For $\varphi \in
(\Vbigalpha^{[\alpha]})^*$:
$(\Id{V_\alpha} \otimes \varphi)(\bv) = \sum_\nu \varphi(\bU_\nu) u_\nu
\in \spann\{u_\nu\} = \Umin{\alpha}{\bv}$. Conversely, by taking
$\varphi = \varphi_\nu$ dual to a basis of $\Umin{D\setminus\{\alpha\}}{\bv}$,
one recovers a basis of $\Umin{\alpha}{\bv}$, so equality holds.
\end{proof}

\begin{remark}
\label{rem:min_replacement}
In Theorem~\ref{thm:min_contraction}, one can replace
$(\Vbigalpha^{[\alpha]})^*$ by $(\Umin{D\setminus\{\alpha\}}{\bv})^*$
without changing the span, since $\bv \in V_\alpha \atens
\Umin{D\setminus\{\alpha\}}{\bv}$. This gives the tighter formula
\begin{equation}
\label{eq:min_contraction_minimal}
    \Umin{\alpha}{\bv}
    = \spann\bigl\{
        (\Id{V_\alpha} \otimes \varphi^{[\alpha]})(\bv) :
        \varphi^{[\alpha]} \in (\Umin{D\setminus\{\alpha\}}{\bv})^*
      \bigr\},
\end{equation}
which is often more useful in practice.
\end{remark}

%------------------------------------------------------------
\section{Lattice properties and inclusion under partitions}
\label{sec:min:lattice}
%------------------------------------------------------------

\subsection*{Monotonicity}

\begin{proposition}[Monotonicity of minimal subspaces]
\label{prop:min_monotone}
Let $\bv \in \VD$ and $U_\alpha \in \Grass{r_\alpha}{V_\alpha}$
for each $\alpha \in D$. Then:
\begin{equation*}
    \bv \in {}_a\bigotimes_{\alpha \in D} U_\alpha
    \quad \Longrightarrow \quad
    \Umin{\alpha}{\bv} \subset U_\alpha \quad \forall\, \alpha \in D.
\end{equation*}
\end{proposition}

\begin{proof}
Under $\bv \in {}_a\bigotimes_{\alpha \in D} U_\alpha$, the identification
$\VD \cong V_\alpha \atens \Vbigalpha^{[\alpha]}$ gives
$\bv \in U_\alpha \atens \Vbigalpha^{[\alpha]}$,
so the minimality of $\Umin{\alpha}{\bv}$ (Theorem~\ref{thm:min_existence}(b))
yields $\Umin{\alpha}{\bv} \subset U_\alpha$.
\end{proof}

\subsection*{Inclusion under partition refinement}

The following proposition is the key algebraic fact underlying the
tree-based format: it shows that minimal subspaces at a coarser
level of a tree are contained in tensor products of minimal subspaces
at finer levels.

\begin{proposition}[Inclusion under refinement]
\label{prop:min_inclusion_partition}
Let $\bv \in \VD$ and let $\alpha \subset D$ with $\card{\alpha} \geq 2$.
Let $\mathcal{P}_\alpha$ be a non-trivial partition of $\alpha$. Then
\begin{equation}
\label{eq:min_inclusion}
    \Umin{\alpha}{\bv} \subset
    {}_a\bigotimes_{\beta \in \mathcal{P}_\alpha} \Umin{\beta}{\bv}.
\end{equation}
\end{proposition}

\begin{proof}
From the canonical identification $\VD \cong \Vbigalpha \atens
\Vbigalpha^{[\alpha]}$ and Theorem~\ref{thm:min_contraction},
\begin{equation*}
    \Umin{\alpha}{\bv} =
    \spann\bigl\{
        (\Id{\Vbigalpha} \otimes \varphi^{[\alpha]})(\bv)
        : \varphi^{[\alpha]} \in (\Umin{\ac}{\bv})^*
    \bigr\}.
\end{equation*}
For each such element $w_\alpha = (\Id{\Vbigalpha} \otimes
\varphi^{[\alpha]})(\bv) \in \Vbigalpha$, we apply
Proposition~\ref{prop:min_monotone} to the tensor space
$\Vbigalpha \cong {}_a\bigotimes_{\beta \in \mathcal{P}_\alpha} \Vbigbeta$.
Since $w_\alpha \in \Umin{\alpha}{\bv}$ and
$\Vbigalpha = {}_a\bigotimes_{\beta \in \mathcal{P}_\alpha} \Vbigbeta$,
one has $w_\alpha \in {}_a\bigotimes_{\beta \in \mathcal{P}_\alpha}
\Umin{\beta}{w_\alpha}$. Moreover, from the minimality of
$\Umin{\beta}{\bv}$ and the relation between contractions at different
levels, $\Umin{\beta}{w_\alpha} \subset \Umin{\beta}{\bv}$ for each
$\beta \in \mathcal{P}_\alpha$. Hence
$w_\alpha \in {}_a\bigotimes_{\beta \in \mathcal{P}_\alpha} \Umin{\beta}{\bv}$,
and taking the span over all such $w_\alpha$ gives~\eqref{eq:min_inclusion}.
\end{proof}

\begin{remark}
\label{rem:min_inclusion_tree}
Proposition~\ref{prop:min_inclusion_partition} applied iteratively at
every internal vertex $\alpha \in \TD$ of a dimension partition tree
(Chapter~\ref{chap:trees}) gives the fundamental compatibility of
minimal subspaces with the tree structure:
\begin{equation}
\label{eq:min_inclusion_tree}
    \Umin{\alpha}{\bv} \subset
    {}_a\bigotimes_{\beta \in \Sons{\alpha}} \Umin{\beta}{\bv}
    \quad \forall\, \alpha \in \TD \setminus \Leaves.
\end{equation}
This is the algebraic foundation of the tree-based format in
Chapter~\ref{chap:algTBT}.
\end{remark}

\subsection*{New characterisation of minimal subspaces}

The following result, proved in~\cite{FHN2021} and used in the construction
of tree-based tensor formats, gives a characterisation of $\Umin{\beta}{\bv}$
via contractions of $\Umin{\alpha}{\bv}$ for $\alpha \supset \beta$.

\begin{proposition}
\label{prop:min_char_partition}
Let $\bv \in \VD$, $\alpha \subset D$ with $\card{\alpha} \geq 2$,
and $\mathcal{P}_\alpha$ a non-trivial partition of $\alpha$. Suppose
that $\Vbigalpha$ and $\Vbigbeta$ for $\beta \in \mathcal{P}_\alpha$
are normed spaces. Then for each $\beta \in \mathcal{P}_\alpha$:
\begin{equation}
\label{eq:min_char_partition}
    \Umin{\beta}{\bv}
    = \spann\bigl\{
        (\Id{\Vbigbeta} \otimes \varphi^{(\alpha\setminus\beta)})
        (\bv_\alpha) :
        \bv_\alpha \in \Umin{\alpha}{\bv},\;
        \varphi^{(\alpha\setminus\beta)} \in
        ({}_a\bigotimes_{\gamma \in \mathcal{P}_\alpha \setminus \{\beta\}}
        \Vbigalpha[\gamma])^*
    \bigr\}.
\end{equation}
\end{proposition}

\begin{proof}
Applying Theorem~\ref{thm:min_contraction} twice — first for the
full index set $D$ and then for $\alpha$ — one obtains two nested
contraction characterisations. The formula~\eqref{eq:min_char_partition}
follows by composing these and using that $\bv \in \Vbigalpha \atens
\Umin{\ac}{\bv}$ allows the replacement of the ambient dual by that
of $\Umin{\ac}{\bv}$.
\end{proof}

%------------------------------------------------------------
\section{Admissible ranks and the disjoint decomposition}
\label{sec:min:admissible}
%------------------------------------------------------------

\subsection*{Admissible Tucker ranks}

Recall from Chapter~\ref{chap:algtensor} (Definition~\ref{def:admissible_tucker})
that the set of admissible Tucker ranks is
\begin{equation*}
    \AD{\VD} = \bigl\{
        (r_\alpha(\bv))_{\alpha \in D} \in \NNz^d : \bv \in \VD
    \bigr\}.
\end{equation*}
The following proposition characterises which tuples are admissible.

\begin{proposition}
\label{prop:admissible_characterisation}
A tuple $\fr = (r_\alpha)_{\alpha \in D} \in \NNz^d$ belongs to
$\AD{\VD}$ if and only if there exist subspaces
$U_\alpha \in \Grass{r_\alpha}{V_\alpha}$ and a tensor
$\bv \in \Tucker{\fr}{{}_a\bigotimes_{\alpha \in D} U_\alpha}$,
i.e., a tensor whose $\alpha$-unfolding has rank exactly $r_\alpha$
for all $\alpha \in D$.
\end{proposition}

\begin{proposition}
\label{prop:admissible_bounds}
For $\fr \in \AD{\VD}$:
\begin{enumerate}
\item[\textup{(i)}] $r_\alpha \leq \dim V_\alpha$ for all $\alpha \in D$.
\item[\textup{(ii)}] $r_\alpha = r_{\alpha^c} \coloneqq
    \dim \Umin{D\setminus\{\alpha\}}{\bv}$ for all $\alpha \in D$
    \textup{(}symmetry of marginal ranks\textup{)}.
\item[\textup{(iii)}] $r_\alpha \leq \prod_{\beta \neq \alpha} r_\beta$
    for all $\alpha \in D$.
\end{enumerate}
\end{proposition}

\begin{proof}
Statements~(i) and~(iii) follow from the representation formula
\eqref{eq:tucker_core} and rank properties of multi-index arrays.
Statement~(ii) follows from Proposition~\ref{prop:min_d2} applied
to $V_\alpha \atens \Vbigalpha^{[\alpha]}$.
\end{proof}

\subsection*{The disjoint decomposition}

Theorem~\ref{thm:min_existence} immediately gives the following
structural result.

\begin{theorem}
\label{thm:tucker_decomposition_full}
The algebraic tensor space $\VD$ admits the disjoint decomposition
\begin{equation}
\label{eq:tucker_disjoint_full}
    \VD = \bigsqcup_{\fr \in \AD{\VD}} \Tucker{\fr}{\VD}.
\end{equation}
Each $\Tucker{\fr}{\VD}$ is non-empty for $\fr \in \AD{\VD}$.
The zero tensor $\mathbf{0}$ belongs to $\Tucker{\mathbf{0}}{\VD}$
and each elementary tensor $\bigotimes_\alpha v_\alpha$ (with all
$v_\alpha \neq 0$) belongs to $\Tucker{\mathbf{1}}{\VD}$.
\end{theorem}

\begin{remark}[Topology of the decomposition]
\label{rem:tucker_topology}
While $\Tucker{\fr}{\VD}$ is algebraically well-defined, its
topological structure is non-trivial. The key result —
that each $\Tucker{\fr}{\VD}$ is a $\mathcal{C}^\infty$-Banach
manifold immersed in the ambient Banach tensor space $\VDnorm$ —
is established in Chapter~\ref{chap:tucker_manifold}. The immersion
requires the condition \condINJ\ on the norm $\|\cdot\|_D$
(Definition~\ref{def:crossnorm}).
\end{remark}

\subsection*{Stratification by rank}

\begin{definition}
\label{def:bounded_tucker_rank}
For $\fr \in \AD{\VD}$, the set of tensors with Tucker rank
\emph{bounded by} $\fr$ is
\begin{equation}
\label{eq:tucker_bounded}
    \Tucker{\leq\fr}{\VD}
    \coloneqq
    \bigsqcup_{\substack{\mathfrak{s} \in \AD{\VD} \\ \mathfrak{s} \leq \fr}}
    \Tucker{\mathfrak{s}}{\VD}
    =
    \bigl\{
        \bv \in \VD : r_\alpha(\bv) \leq r_\alpha \;\forall\, \alpha \in D
    \bigr\},
\end{equation}
where $\fs \leq \fr$ means $s_\alpha \leq r_\alpha$ for all $\alpha \in D$.
\end{definition}

The set $\Tucker{\leq\fr}{\VD}$ is not a linear subspace (it is not
closed under addition) but it is closed under scalar multiplication.
The key property for applications — that best approximations from
$\Tucker{\leq\fr}{\VD}$ exist in Banach tensor spaces — is proved in
Section~\ref{sec:min:bestapprox}.

%------------------------------------------------------------
\section{Best approximations in Banach tensor spaces}
\label{sec:min:bestapprox}
%------------------------------------------------------------

In numerical applications, one needs to approximate a given tensor
$\bv \in \VDnorm$ by a tensor of bounded rank. The existence of such
a best approximation is non-trivial: the set $\Tucker{\leq\fr}{\VD}$
is not convex, so general theorems about proximinal sets do not apply.
The proof uses the weak closedness of $\Tucker{\leq\fr}{\VD}$.

\subsection*{Crossnorms and the injective norm condition}

\begin{definition}
\label{def:crossnorm}
Let $(V_\alpha, \|\cdot\|_\alpha)$ be normed spaces and $\|\cdot\|_D$
a norm on $\VD$. We say $\|\cdot\|_D$ is a \emph{crossnorm} if
\begin{equation}
\label{eq:crossnorm}
    \Bigl\|\bigotimes_{\alpha \in D} v_\alpha\Bigr\|_D
    = \prod_{\alpha \in D} \|v_\alpha\|_\alpha
    \quad \forall\, v_\alpha \in V_\alpha.
\end{equation}
The \emph{injective norm} $\injnorm$ on $\VD$ is defined by
\begin{equation*}
    \|\bv\|_\varepsilon
    \coloneqq \sup\bigl\{
        |(\bigotimes_{\alpha \in D} \varphi_\alpha)(\bv)| :
        \varphi_\alpha \in \dual{V_\alpha},\;
        \|\varphi_\alpha\|_{\dual{V_\alpha}} \leq 1
    \bigr\}.
\end{equation*}
We say $\|\cdot\|_D$ satisfies condition \condINJ\ if
$\|\bv\|_D \geq \|\bv\|_\varepsilon$ for all $\bv \in \VD$.
\end{definition}

A crossnorm satisfies \condINJ\ if and only if it is a
\emph{reasonable crossnorm} in the sense that its dual norm is also
a crossnorm. The injective norm is the smallest reasonable crossnorm
(see \cite{Hackbusch2019}, Chapter~4).

\begin{remark}
\label{rem:crossnorm_examples}
The $L^p$ norm on $L^p(\prod_\alpha I_\alpha)$ is a reasonable
crossnorm for $1 \leq p \leq \infty$
(see \cite{Hackbusch2019}, Example~4.72).
The Sobolev norm $\|\cdot\|_{(0,1),p}$ on
$H^{1,p}(I_1) \atens H^{1,p}(I_2)$ is a crossnorm
(Example~4.42 in \cite{Hackbusch2019}), hence satisfies \condINJ.
\end{remark}

\subsection*{Weak closedness of bounded-rank sets}

\begin{theorem}[Weak closedness, {\cite[Theorem~5.1]{FH2012}}]
\label{thm:weak_closed}
Let $(V_\alpha, \|\cdot\|_\alpha)$ be normed spaces and let
$\|\cdot\|_D$ be a norm on $\VD$ satisfying condition \condINJ.
Then the set $\Tucker{\leq\fr}{\VD}$ is weakly closed in $\VDnorm$.
\end{theorem}

\begin{proof}[Sketch]
Let $(\bv_n) \subset \Tucker{\leq\fr}{\VD}$ converge weakly to
$\bv \in \VDnorm$. We must show $\bv \in \Tucker{\leq\fr}{\VD}$,
i.e., $r_\alpha(\bv) \leq r_\alpha$ for all $\alpha \in D$.

Fix $\alpha \in D$ and consider the $\alpha$-unfolding map
$\Mmat{\alpha}: \VDnorm \to \cL(\dual{(V_\alpha)}, \Vbigalpha^{[\alpha]})$
defined by $\Mmat{\alpha}(\bv)(\varphi) = (\varphi \otimes \Id{})(\bv)$.
Since $r_\alpha(\bv_n) \leq r_\alpha$, we have
$\rank \Mmat{\alpha}(\bv_n) \leq r_\alpha$ for all $n$. The condition
\condINJ\ implies that $\Mmat{\alpha}$ is continuous for the weak
topology on $\VDnorm$ and the weak operator topology on
$\cL(\dual{(V_\alpha)}, \Vbigalpha^{[\alpha]})$. Since the set of
operators of rank $\leq r_\alpha$ is closed in the weak operator topology
(it is a closed variety), the limit $\Mmat{\alpha}(\bv)$ has rank
$\leq r_\alpha$, i.e., $r_\alpha(\bv) \leq r_\alpha$.
\end{proof}

\subsection*{Existence of best approximations}

\begin{theorem}[Best approximation, {\cite[Theorem~5.3]{FH2012}}]
\label{thm:best_approx_tucker}
Let $(V_\alpha, \|\cdot\|_\alpha)$ be normed spaces, let $\|\cdot\|_D$
be a norm on $\VD$ satisfying \condINJ, and assume that $\VDnorm$ is
reflexive. Then for every $\bu \in \VDnorm$ and $\fr \in \AD{\VD}$,
there exists $\bv_* \in \Tucker{\leq\fr}{\VDnorm}$ such that
\begin{equation}
\label{eq:best_approx}
    \|\bu - \bv_*\|_D = \inf_{\bv \in \Tucker{\leq\fr}{\VDnorm}}
    \|\bu - \bv\|_D.
\end{equation}
\end{theorem}

\begin{proof}
Since $\VDnorm$ is reflexive, any bounded sequence has a weakly
convergent subsequence (Eberlein--Šmulian theorem). Let
$(\bv_n) \subset \Tucker{\leq\fr}{\VDnorm}$ be a minimising sequence
for the infimum in~\eqref{eq:best_approx}. The sequence $(\bv_n)$ is
bounded (the infimum is attained or approached by a bounded set), so
by reflexivity there is a weakly convergent subsequence
$\bv_{n_k} \rightharpoonup \bv_*$. By
Theorem~\ref{thm:weak_closed}, $\bv_* \in \Tucker{\leq\fr}{\VDnorm}$.
Since the norm is weakly lower semicontinuous,
$\|\bu - \bv_*\|_D \leq \liminf_k \|\bu - \bv_{n_k}\|_D = \inf$,
so $\bv_*$ is a best approximation.
\end{proof}

\begin{corollary}
\label{cor:best_approx_lp}
For $I_\alpha \subset \RR$ ($\alpha \in D$) and $1 < p < \infty$,
the Banach tensor space $L^p(\prod_\alpha I_\alpha)$ is reflexive and
its norm satisfies \condINJ\ (Example~\ref{rem:crossnorm_examples}).
Hence every $f \in L^p(\prod_\alpha I_\alpha)$ has a best
Tucker approximation of any prescribed rank.
\end{corollary}

\begin{corollary}
\label{cor:best_approx_sobolev}
For $1 < p < \infty$, $N \geq 1$, and bounded intervals $I_\alpha$:
\begin{enumerate}
\item[\textup{(i)}] $H^{N,p}(\prod_\alpha I_\alpha)$ with the standard
    Sobolev norm is reflexive and satisfies \condINJ.
\item[\textup{(ii)}] The space $H^{1,p}(I_1) \otimes_{\|\cdot\|_{(0,1),p}}
    H^{1,p}(I_2)$ with the partial-derivative crossnorm satisfies \condINJ.
\end{enumerate}
In both cases, best Tucker approximations of any rank exist.
\end{corollary}

\begin{example}
\label{ex:min:sobolev}
Consider $d=2$, $V_1 = H^1(0,1)$, $V_2 = H^1(0,1)$, and the
ambient space $V_1 \otimes_{H^1} V_2$ with the norm
$\|f\|_{H^1}^2 = \|f\|_{L^2}^2 + \|\partial_{x_2} f\|_{L^2}^2$.
This Banach tensor space is reflexive and satisfies \condINJ, so every
$f \in V_1 \otimes_{H^1} V_2$ has a best approximation of
rank $\leq r$ for any $r \geq 1$. This setting arises in parametric PDEs
where $x_1$ is the spatial variable and $x_2$ is the stochastic parameter.
\end{example}

\subsection*{Extension to intersections}

In many applications (including the tree-based format), the ambient
space is an intersection of Banach tensor spaces, each with its own norm.

\begin{theorem}[Best approximation in intersections, {\cite[Theorem~5.6]{FH2012}}]
\label{thm:best_approx_intersection}
Let $\VDnorm^{(1)}, \ldots, \VDnorm^{(k)}$ be reflexive Banach tensor
spaces over the same algebraic tensor space $\VD$, each with a norm
satisfying \condINJ. Equip the intersection
$Z = \VDnorm^{(1)} \cap \cdots \cap \VDnorm^{(k)}$ with the norm
$\|\bv\|_Z = \max_{i} \|\bv\|_i$. If $Z$ is reflexive (which holds,
e.g., when $k=2$ and both norms are reflexive), then every $\bu \in Z$
has a best approximation in $\Tucker{\leq\fr}{Z}$.
\end{theorem}

%------------------------------------------------------------
\section{Worked examples}
\label{sec:min:examples}
%------------------------------------------------------------

This section illustrates the abstract theory through two concrete
examples that arise naturally in scientific computing.

\subsection*{Example 1: A Gaussian kernel in $L^2([0,1]^2)$}

\begin{example}[Gaussian kernel]
\label{ex:min:gaussian}
Let $V_1 = V_2 = L^2([0,1])$ and consider the function
\begin{equation}
\label{eq:gaussian_kernel}
    f(x_1, x_2) \coloneqq e^{-\gamma(x_1 - x_2)^2},
    \quad (x_1, x_2) \in [0,1]^2,\quad \gamma > 0.
\end{equation}
We identify $f$ with a Hilbert--Schmidt operator
$T_f: L^2([0,1]) \to L^2([0,1])$ via
$(T_f g)(x_1) = \int_0^1 f(x_1, x_2) g(x_2)\,dx_2$.

\paragraph{Minimal subspaces.}
Since $T_f$ is a compact self-adjoint positive operator
(the kernel is symmetric and positive definite), its spectral
decomposition is
\begin{equation*}
    T_f = \sum_{k=1}^\infty \lambda_k\, \phi_k \otimes \phi_k,
    \quad \lambda_1 \geq \lambda_2 \geq \cdots > 0, \quad
    \sum_k \lambda_k < \infty,
\end{equation*}
with $\{\phi_k\}_{k \geq 1}$ an orthonormal basis of $L^2([0,1])$
consisting of eigenfunctions of $T_f$. In the tensor product notation,
\begin{equation*}
    f = \sum_{k=1}^\infty \lambda_k\, \phi_k \otimes \phi_k
    \in L^2([0,1]) \otimes_H L^2([0,1]),
\end{equation*}
so $r_1(f) = r_2(f) = \infty$ (the kernel has infinite rank). For any
$r \in \NN$, however, the minimal subspaces of the rank-$r$ truncation
\begin{equation*}
    f_r \coloneqq \sum_{k=1}^r \lambda_k\, \phi_k \otimes \phi_k
\end{equation*}
are $\Umin{1}{f_r} = \Umin{2}{f_r} = \spann\{\phi_1, \ldots, \phi_r\}$,
and $f_r \in \Tucker{(r,r)}{L^2([0,1])^{\otimes 2}}$.

\paragraph{Best approximation.}
By the Eckart--Young theorem (the $d=2$ case of
Theorem~\ref{thm:best_approx_tucker}), $f_r$ is the
best approximation of $f$ of Tucker rank $(r,r)$ in the Hilbert tensor
norm:
\begin{equation*}
    \|f - f_r\|_{L^2([0,1]^2)}^2
    = \sum_{k > r} \lambda_k^2
    = \inf_{\bg \in \Tucker{\leq(r,r)}{L^2}} \|f - \bg\|^2.
\end{equation*}

\paragraph{Decay of eigenvalues.}
For $\gamma > 0$, the eigenvalues of the Gaussian kernel
decay super-algebraically: $\lambda_k \sim C e^{-c\sqrt{k}}$
for constants $C, c > 0$ depending on $\gamma$. This rapid decay
means that $f_r$ approximates $f$ to any prescribed accuracy with
$r$ growing only polylogarithmically with the tolerance. More
precisely, for tolerance $\varepsilon > 0$:
\begin{equation}
\label{eq:gaussian_rank}
    \|f - f_r\|_{L^2}^2 \leq \varepsilon^2
    \quad \text{requires } r = O\bigl(\log(1/\varepsilon)^2\bigr).
\end{equation}
This is the mathematical reason why low-rank approximations work
exceptionally well for functions with smooth integral kernels: the
minimal subspace dimension needed for $\varepsilon$-accuracy grows only
logarithmically with $1/\varepsilon$.
\end{example}

\subsection*{Example 2: Parametric partial differential equations}

The main motivation for the tensor framework of this monograph is the
numerical treatment of high-dimensional PDEs. The following example
shows how the minimal subspace dimension encodes the sensitivity of the
solution with respect to parameters.

\begin{example}[Parametric elliptic PDE]
\label{ex:min:parametric_pde}
Let $D_x = (0,1) \subset \RR$ and $\Gamma = [-1,1]^s \subset \RR^s$
be a parameter space ($s \geq 1$). Consider the parametric family of
elliptic problems: find $u(\cdot, \bm{y}): D_x \to \RR$ such that
\begin{equation}
\label{eq:parametric_pde}
    -\frac{\partial}{\partial x}
    \Bigl[a(x, \bm{y})\,\frac{\partial u}{\partial x}(x, \bm{y})\Bigr]
    = f(x),
    \quad x \in D_x,\; \bm{y} = (y_1, \ldots, y_s) \in \Gamma,
\end{equation}
with homogeneous Dirichlet boundary conditions $u(0,\bm{y}) =
u(1,\bm{y}) = 0$. The diffusion coefficient has the affine form
\begin{equation*}
    a(x, \bm{y}) = a_0(x) + \sum_{j=1}^s y_j a_j(x),
\end{equation*}
where $a_0, a_1, \ldots, a_s \in L^\infty(D_x)$ with
$a(x, \bm{y}) \geq a_{\min} > 0$ uniformly.

\paragraph{Tensor product setting.}
The solution map $\bm{y} \mapsto u(\cdot, \bm{y})$ defines an element
\begin{equation*}
    \bm{u} \in \mathcal{V} \coloneqq
    H_0^1(D_x) \otimes L^2(\Gamma, \pi),
\end{equation*}
where $\pi$ is the uniform measure on $\Gamma$ and the tensor product
norm satisfies \condINJ\ (Corollary~\ref{cor:best_approx_sobolev}).

\paragraph{Minimal subspaces and parameter sensitivity.}
The $\alpha$-rank $r_\alpha(\bm{u})$ for $\alpha = \{x\}$ (the physical
variable) equals the effective rank of the solution manifold
$\{u(\cdot, \bm{y}) : \bm{y} \in \Gamma\} \subset H_0^1(D_x)$:
\begin{equation*}
    \Umin{x}{\bm{u}}
    = \overline{\spann\{u(\cdot, \bm{y}) : \bm{y} \in \Gamma\}}
    \subset H_0^1(D_x).
\end{equation*}
If $a_0$ dominates (i.e., the perturbations $y_j a_j$ are small),
the solution $u(\cdot, \bm{y})$ varies little with $\bm{y}$, and
$r_x(\bm{u})$ is small. More precisely, the Kolmogorov $n$-width
\begin{equation*}
    d_n\bigl(\{u(\cdot,\bm{y}) : \bm{y}\in\Gamma\},\,H_0^1(D_x)\bigr)
    \coloneqq \inf_{\substack{W \subset H_0^1(D_x) \\ \dim W = n}}
    \sup_{\bm{y} \in \Gamma}
    \inf_{w \in W} \|u(\cdot,\bm{y}) - w\|_{H_0^1}
\end{equation*}
decays at a rate determined by the smoothness of $a$ in the parameters
$\bm{y}$. When $a$ is analytic in $\bm{y}$ (which holds when the
$a_j$ are bounded away from zero), the $n$-width decays exponentially:
$d_n \leq C e^{-cn^{1/s}}$ for constants $C, c > 0$
(see~\cite{Nouy2017}).

\paragraph{Best approximation from Tucker format.}
The best Tucker-$(r,r)$ approximation of $\bm{u}$ in $\mathcal{V}$
is
\begin{equation*}
    \bm{u}_r = \sum_{k=1}^r \lambda_k\, v_k^{(x)} \otimes v_k^{(\bm{y})},
    \quad
    \|\bm{u} - \bm{u}_r\|_{\mathcal{V}}^2 = \sum_{k > r} \lambda_k^2,
\end{equation*}
where $\lambda_k$ are the singular values of the solution operator
$\bm{y} \mapsto u(\cdot, \bm{y})$ as a Hilbert--Schmidt operator. The
existence of this best approximation is guaranteed by
Corollary~\ref{cor:best_approx_sobolev}. The spatial basis functions
$v_k^{(x)} \in H_0^1(D_x)$ are the principal components of the
solution manifold, and $v_k^{(\bm{y})} \in L^2(\Gamma)$ encode the
parameter dependence. This decomposition is the starting point for
the \emph{proper orthogonal decomposition} (POD) and reduced basis
methods used in uncertainty quantification; see~\cite{Nouy2017}.
\end{example}

%------------------------------------------------------------
\section{Historical notes and comparison with the literature}
\label{sec:min:history}
%------------------------------------------------------------

The notion of minimal subspace was introduced in~\cite{FH2012}.
It generalises to infinite dimensions the classical concept of matrix
rank: for $d=2$ and finite-dimensional $V_\alpha$, the $\alpha$-rank
$r_\alpha(\bv) = \rank \Mmat{\alpha}(\bv)$ coincides with the classical
matrix rank of the $\alpha$-unfolding, introduced by
Hitchcock~\cite{Hitchcock1927} in 1927. The higher-order version for
finite-dimensional spaces was studied by Tucker~\cite{Tucker1966},
De Lathauwer et al.~\cite{DeLathauwer2000}, and Kolda--Bader~\cite{KoldaBader2009}.

The infinite-dimensional theory of~\cite{FH2012} is motivated by the
numerical treatment of high-dimensional PDEs, where $V_\alpha = L^2(\RR^3)$
or Sobolev spaces, and the ambient tensor space is infinite-dimensional.
The existence of best approximations (Theorem~\ref{thm:best_approx_tucker})
removes a fundamental obstacle in the convergence analysis of alternating
least-squares and related iterative methods.

The weak closedness argument of Theorem~\ref{thm:weak_closed} is
modelled on Uschmajew's result for the case of the spectral norm
on finite-dimensional spaces~\cite{Uschmajew2012}. The innovation
of~\cite{FH2012} is the extension to general reflexive Banach spaces
under the crossnorm condition \condINJ, which is the natural condition
for tensor product norms.

The connection between minimal subspaces and the geometry of Tucker
manifolds is developed in Chapter~\ref{chap:tucker_manifold}, where
the rank map $\varrho: \VD \to \prod_\alpha \Grass{r_\alpha}{V_\alpha}$
defined by $\varrho(\bv) = (\Umin{\alpha}{\bv})_{\alpha \in D}$ plays
a central role in the construction of local charts.

